\lecture{3}{9. Februar 2026}{The Cauchy Theorem \& Principal Stresses}

\begin{problem}{2.1}
  Find the components of the traction on a plane defined by $n_1 = 1 / \sqrt{2}$, $n_2 = 1 / \sqrt{2}$, and $n_3 = 0$ for the following states of stress:
  \begin{enumerate}[label=(\alph*)]
    \item $\sigma_{11} = \sigma, \sigma_{12} = 0, \sigma_{13} = 0$, $\sigma_{21} = \sigma$, $\sigma_{22} = 0$, $\sigma_{33} = \sigma$.
    \item $\sigma_{11} = \sigma$, $\sigma_{12} = \sigma$, $\sigma_{13} = 0$, $\sigma_{22} = \sigma$, $\sigma_{23} = 0$, $\sigma_{33} = 0$.
  \end{enumerate}
\end{problem}

\begin{derivation}
  \begin{enumerate}[label=(\alph*)]
    \item The stress state in (a) is
      \begin{equation*}
        \sigma = \begin{bmatrix}
          \sigma & 0 & 0\\
        \sigma & 0 & \sigma_{23}\\
        \sigma_{31} & \sigma_{32} & \sigma\\
        \end{bmatrix}
      .\end{equation*}
      The Cauchy Theorem states that
      \begin{equation*}
        T_i = \sigma_{ji} n_j \implies T = \sigma \cdot n
      .\end{equation*}
      Our plane is:
      \begin{equation*}
              n = \begin{bmatrix}
                \frac{1}{\sqrt{2}}\\
              \frac{1}{\sqrt{2}}\\
              0\\
              \end{bmatrix}
      .\end{equation*}
      Hence,
      \begin{equation*}
        T = \begin{bmatrix}
          \sigma & 0 & 0\\
        \sigma & 0 & \sigma_{23}\\
        \sigma_{31} & \sigma_{32} & \sigma\\
        \end{bmatrix} \cdot \begin{bmatrix}
          \frac{1}{\sqrt{2}}\\
        \frac{1}{\sqrt{2}}\\
        0\\
        \end{bmatrix} = \left(
\begin{array}{c}
 \frac{\sigma }{\sqrt{2}} \\
 \frac{\sigma }{\sqrt{2}} \\
 \frac{\sigma _{31}}{\sqrt{2}}+\frac{\sigma _{32}}{\sqrt{2}} \\
\end{array}
\right)
      .\end{equation*}
    \item For $(b)$ our stress state is:
      \begin{equation*}
        \sigma = \begin{bmatrix}
          \sigma & \sigma & 0\\
        \sigma_{21} & \sigma & 0\\
        \sigma_{31} & \sigma_{32} & 0\\
        \end{bmatrix}
      ,\end{equation*}
      so
      \begin{equation*}
        T = \begin{bmatrix}
          \sigma & \sigma & 0\\
        \sigma_{21} & \sigma & 0\\
        \sigma_{31} & \sigma_{32} & 0\\
        \end{bmatrix} \cdot \begin{bmatrix}
          \frac{1}{\sqrt{2}}\\
        \frac{1}{\sqrt{2}}\\
        0\\
        \end{bmatrix} = \left(
\begin{array}{c}
 \sqrt{2} \sigma  \\
 \frac{\sigma }{\sqrt{2}}+\frac{\sigma _{21}}{\sqrt{2}} \\
 \frac{\sigma _{31}}{\sqrt{2}}+\frac{\sigma _{32}}{\sqrt{2}} \\
\end{array}
\right)
      .\end{equation*}
  \end{enumerate}
\end{derivation}

\finalresult{
  \begin{aligned}
    &(\mathrm{a}) & T &= \begin{bmatrix}
      \frac{\sigma}{\sqrt{2}}\\
    \frac{\sigma}{\sqrt{2}}\\
    \frac{\sigma_{31}}{\sqrt{2}} + \frac{\sigma_{32}}{\sqrt{2}}\\
    \end{bmatrix} \\
    &(\mathrm{b}) & T &= \begin{bmatrix}
      \sqrt{2}\sigma\\
    \frac{\sigma}{\sqrt{2}} + \frac{\sigma_{21}}{\sqrt{2}}\\
    \frac{\sigma_{31}}{\sqrt{2}} + \frac{\sigma_{32}}{\sqrt{2}}\\
    \end{bmatrix}
  .\end{aligned}
}

\begin{problem}{2.6}
  At a given point in a body,
  \begin{align*}
    \sigma_{11} = \sigma_{22} = \sigma_{33} &= -p, \\
    \sigma_{12} = \sigma_{23} = \sigma_{31} &= 0
  .\end{align*}
  Show that the normal stresses are equal to $-p$ and that the shearing stresses vanish for any other Cartesian coordinate system.
\end{problem}

\begin{derivation}
  Here, our stress tensor is
  \begin{equation*}
    \Vec{\sigma} = \begin{bmatrix}
      -p & 0 & 0\\
    0 & -p & 0\\
    0 & 0 & -p\\
    \end{bmatrix} = -p \Vec{I}
  .\end{equation*}
  Using the Cauchy theorem $\Vec{T} = \Vec{\sigma} \cdot \Vec{n}$ we get:
  \begin{equation*}
    \Vec{T} = (-p \Vec{I}) \cdot \Vec{n} = -p \Vec{n}
  .\end{equation*}
  Hence, $\Vec{T}$ is parallel to $\Vec{n}$ meaning the traction does not have any components tangent to the plane and hence the stress is entirely normal with size $-p$.
\end{derivation}
